\documentclass[11pt]{article}
\usepackage{amsmath, amssymb, amsthm}

\title{Corrected Structural Placement of the Five Threshold Domains:\\
AI Interpretability Altitudes and Soft-Manifold Embeddings}
\author{Borealis S. Hedling}
\date{}

\begin{document}
\maketitle

\section*{Abstract}
The original Axiom Quant ``Five Threshold Domains'' infographic presents five
independent phase-transition systems as if they share a universal mathematical
pipeline. This document provides a corrected, altitude-clean interpretation by
embedding each domain into an 80-dimensional soft Riemannian manifold and
mapping each threshold phenomenon to its appropriate AI interpretability
altitude. The result is a coherent geometric and interpretability-based
framework that avoids false equivalences while preserving conceptual unity
through invariant extraction.

\section{The Parent Geometry: Soft 80D Riemannian Manifold}
Let $(\mathcal{M}^{80}, g)$ be a soft Riemannian manifold with reversible
curvature envelopes and holonomy fields. Softness ensures that perturbations
do not induce topological collapse:
\[
\|\Delta g\| < \varepsilon \quad \Rightarrow \quad \text{no singularity forms}.
\]
Each threshold domain is embedded into a distinct region of $\mathcal{M}^{80}$,
with adjacency defined by curvature gradients rather than algebraic identity.

\section{AI Interpretability Altitudes}
We adopt the altitude model:
\[
\text{A3: Mechanistic},\quad
\text{A4: Feature Circuits},\quad
\text{A5: Latent Semantics},\quad
\text{A7--A10: Geometric Manifolds}.
\]
Each threshold domain is assigned to the altitude where its invariants become
interpretable.

\section{Domain 1: Erd\H{o}s Graph Threshold}
Edge probability $p$ induces a connectivity phase transition. We embed the
graph family into a region $\mathcal{R}_{\mathrm{graph}} \subset \mathcal{M}^{80}$
with curvature proportional to component fragmentation:
\[
K(p) \propto \text{component entropy}.
\]
Interpretability altitude: A3 (mechanistic graph structure).

\section{Domain 2: BitNet Ternary Quantization}
The scale factor $\gamma$ controls the transition from continuous weights to
ternary $\{-1,0,+1\}$. We embed quantization surfaces into
$\mathcal{R}_{\mathrm{ternary}}$ with curvature slack allowing discrete jumps
without metric discontinuity. Holonomy tracks rounding drift:
\[
\mathcal{H}(\gamma) = \oint_\gamma \omega_{\mathrm{quant}}.
\]
Interpretability altitude: A4--A5 (feature circuits and latent semantics).

\section{Domain 3: Markowitz Portfolio Threshold}
The Lagrange multiplier $\lambda$ determines asset activation. We embed
efficient-frontier curves into $\mathcal{R}_{\mathrm{portfolio}}$ where convex
geometry is preserved:
\[
\nabla^2 f(x) \succ 0 \quad \Rightarrow \quad \text{stable curvature basin}.
\]
Interpretability altitude: A5 (latent semantic optimization geometry).

\section{Domain 4: Black--Scholes Volatility Threshold}
Volatility $\sigma$ induces transitions between OTM, ATM, and ITM regimes.
We embed option surfaces into $\mathcal{R}_{\mathrm{vol}}$ with positive
curvature corresponding to the canonical form:
\[
\Omega_{\mathrm{BS}} = N(d_2)\, d\sigma.
\]
Interpretability altitude: A7 (spectral bundles).

\section{Domain 5: Rényi Dyadic Map Threshold}
The parameter $\beta$ controls information retention vs.\ chaotic shift. We
embed dyadic dynamics into $\mathcal{R}_{\mathrm{dyadic}}$ with holonomy loops
representing entropy drift:
\[
\Delta H = \beta \cdot \text{shift rate}.
\]
Interpretability altitude: A8--A10 (neutral and soft manifolds).

\section{Unified Structural Correspondence}
We define a correspondence map:
\[
\Xi : \{\text{Graph},\ \text{Ternary},\ \text{Portfolio},\ 
\text{Volatility},\ \text{Dyadic}\}
\to \mathcal{S},
\]
where $\mathcal{S}$ is the invariant-extraction space:
\[
\mathcal{S} = \{\text{connectivity},\ \text{quantization},\ 
\text{convexity},\ \text{canonical form},\ \text{entropy}\}.
\]
The manifold embedding ensures that these correspondences arise from geometric
adjacency rather than forced algebraic equivalence.

\section{Conclusion}
The corrected framework places each threshold domain in its appropriate
interpretability altitude and embeds each into a distinct region of an 80D soft
Riemannian manifold. This resolves the inconsistencies of the original
infographic and provides a coherent geometric foundation for cross-domain
analysis without violating metric constraints or relying on softmax-style
blending.

\end{document}
